\documentclass{article}
\usepackage[margin=1in]{geometry}
\usepackage[T1]{fontenc}
\usepackage{lmodern}
\usepackage{hyperref}
\usepackage{listings}
\lstset{
	language=Python,
	basicstyle=\ttfamily\small,
	keywordstyle=\bfseries,
	stringstyle=\itshape,
	commentstyle=\itshape,
	numbers=none,
	frame=single,
	breaklines=true,
	showstringspaces=false
}
\begin{document}
	
	\section*{Setting Up \texttt{observer-agent-mvp} on Windows}
	
	This guide shows how to get the project running end‐to‐end:
	
	\begin{enumerate}
		\item Create and activate a virtual environment.
		\item Patch \texttt{bootstrap.py} for Windows.
		\item Install dependencies.
		\item Adjust \texttt{auto.py} invocations to use the venv interpreter.
		\item Ensure the \texttt{recall} module finds \texttt{config.py}.
		\item Fix console‐script entry points and add \texttt{main()}.
	\end{enumerate}
	
	\subsection*{1. Virtual Environment}
	
	\begin{lstlisting}[frame=none]
		python -m venv .venv
		.\.venv\Scripts\Activate.bat
	\end{lstlisting}
	
	\subsection*{2. Patch \texttt{bootstrap.py}}
	
	In the \texttt{setup(…)} call, change:
	\begin{lstlisting}[frame=none]
		packages=find_packages(),
		entry_points={
			'console_scripts': [
			'screen-recorder=recording:main',
			],
		},
	\end{lstlisting}
	to:
	\begin{lstlisting}[frame=none]
		packages=find_packages() + ['scripts'],
		entry_points={
			'console_scripts': [
			'screen-recorder=scripts.recording:main',
			],
		},
	\end{lstlisting}
	Also, create an empty file:
	\begin{lstlisting}[frame=none]
		scripts/__init__.py
	\end{lstlisting}
	
	\subsection*{3. Install Requirements}
	
	\begin{lstlisting}[frame=none]
		pip install -r requirements.txt
	\end{lstlisting}
	
	\subsection*{4. Update \texttt{auto.py}}
	
	At the top, add \texttt{import sys}.  
	Replace every
	\begin{lstlisting}[frame=none]
		subprocess.run(['python3', script_path])
	\end{lstlisting}
	with
	\begin{lstlisting}[frame=none]
		subprocess.run([sys.executable, script_path], check=True)
	\end{lstlisting}
	and similarly adjust any
	\begin{lstlisting}[frame=none]
		cmd = ['python3', script_path, ...]
	\end{lstlisting}
	to
	\begin{lstlisting}[frame=none]
		cmd = [sys.executable, script_path, ...]
	\end{lstlisting}
	
	\subsection*{5. Running the \texttt{recall} Module}
	
	Always invoke from the project root so that \texttt{config.py} is on the path:
	\begin{lstlisting}[frame=none]
		python -m scripts.recall
	\end{lstlisting}
	
	\subsection*{6. Bootstrap \& Recording Fix}
	
	After installing via:
	\begin{lstlisting}[frame=none]
		python bootstrap.py install
	\end{lstlisting}
	you will be prompted:
	\begin{lstlisting}[frame=none]
		Would you like to run the screen recorder now? (y/N):
	\end{lstlisting}
	To have \texttt{scripts/recording.py} export a \texttt{main()} for the console‐script, add at the bottom of \texttt{recording.py}:
	\begin{lstlisting}
		def main():
		"""Console-script entry point."""
		start_screen_recording()
		
		if __name__ == "__main__":
		main()
	\end{lstlisting}
	
	Reinstall to pick up the change:
	\begin{lstlisting}[frame=none]
		python bootstrap.py install
	\end{lstlisting}
	
	Now:
	\begin{lstlisting}[frame=none]
		screen-recorder
	\end{lstlisting}
	will launch the recorder immediately after you press \texttt{y}.
	
\end{document}
